\documentclass{article}
\setlength{\oddsidemargin}{-0.3in}
\setlength{\evensidemargin}{-0.3in}
\setlength{\textwidth}{6.5in}
\setlength{\topmargin}{-.7in}
\setlength{\textheight}{9.6in}
\usepackage{amsmath,amsfonts,amssymb,graphicx,fancyhdr,latexsym,float,amsthm,bm,graphicx,mathtools,txfonts,enumitem,listings,verbatim,color}
\usepackage{epsfig}
\usepackage[ruled,vlined, linesnumbered]{algorithm2e}
\usepackage[authoryear,round]{natbib}
\usepackage[
  colorlinks=true,
  urlcolor=blue,
  linkcolor=black,
  citecolor=black
]{hyperref}
\urlstyle{same}
\newcommand{\R}{\mathbb{R}}
\newcommand{\argmax}{\operatornamewithlimits{argmax}}
\newcommand{\argmin}{\operatornamewithlimits{argmin}}
\newcommand{\bS}{\mathcal{S}}

\DeclareMathOperator{\tr}{tr}

\newcounter{chunk}
\parskip=11pt
\parindent=0mm

\pagestyle{fancy}
\lhead{Final Project Proposal}
\rhead{\textit{Group 5}}

\title{Brain Tumor Classification from MRI: A Comparative Study of Convolutional Neural Networks and Vision Transformers}
\author{Enci Zhang, Lufei Wang, Tony Hou, Jessica Nickell}
\date{} %Please don't insert a date

\begin{document}
\maketitle
\section*{Problem Description}
Automated classification of two-dimensional brain MRI images is a challenging image-recognition task because different tumor categories can share visual
characteristics, while images vary in appearance and quality. Convolutional neural networks (CNNs) and vision transformers offer different approaches to
extracting image features. Evaluating these approaches requires considering both classification performance and computational cost, since improvements
in predictive performance may involve larger models or longer training and inference times.

This project investigates the classification of brain MRI images into four categories: glioma, meningioma, pituitary tumor, and no tumor. We will use
version 2 of the publicly available \href{https://www.kaggle.com/datasets/masoudnickparvar/brain-tumor-mri-dataset}{Brain Tumor MRI Dataset}, which contains 
7,200 images, to compare selected ImageNet-1K-pretrained CNN and vision transformer models under a shared experimental protocol. Our central question is how these 
models differ in classification performance, model size, training time, and inference latency when evaluated using the same audited data partitions. The study will
provide a reproducible comparison on this image benchmark. Because patient identifiers are unavailable, its findings will not establish performance on
independent patients or clinical populations.

\section*{Related Work}
Convolutional neural networks (CNNs) are widely used for image classification due to their capacity to learn hierarchical visual representations.
Recent studies have shown that deep learning has achieved strong performance in brain tumor classification using MRI images. 
\citet{mzoughi2020} proposed a multi-scale 3D convolutional neural network that uses volumetric MRI data to classify glioma tumors. 
Their work demonstrated that 3D CNNs can effectively capture spatial information across MRI slices, but they also require more computational 
resources. 

\citet{he2016} introduced residual learning,
which facilitates the optimization of deep networks through shortcut connections. This framework provides the foundation for ResNet-50,
selected in the present study as an established CNN baseline. \citet{tan2019} proposed EfficientNet, a family of architectures developed by
balancing network depth, width, and input resolution. EfficientNet-B0, one model of the family, is introduced to examine the trade-off between 
classification performanceand computational cost within a compact convolutional architecture.

Transformer architectures offer another approach to image classification by using attention to model relationships between image patches.
\citet{touvron2021} introduced Data-Efficient Image Transformers (DeiT), which shows competitive image classification performance through
training on ImageNet without additional external training data. Pretrained DeiT-Small is included to investigate performance of a vision
Transformer as an alternative to convolutional feature extraction for brain MRI classification tasks.

\citet{liu2021} proposed Swin Transformer, which constructs hierarchical representations using self-attention within local windows.
The shifted-window mechanism enables information exchange across window boundaries, while the hierarchical architecture supports feature
representation at multiple spatial scales. Swin-Tiny is included to examine whether this attention structure offers advantages over the global
attention employed by DeiT-Small for the selected classification task.

\section*{Proposed Work}
\setlength{\emergencystretch}{2em}
We will compare ImageNet-1K-pretrained ResNet-50, \mbox{EfficientNet-B0}, DeiT-Small, and Swin-Tiny in classification performance and computational efficiency, 
using the standard, non-distilled DeiT-Small configuration. After replacing each model's classification head with a four-output layer, we will fine-tune all
parameters. For a consistent comparison, all models will share identical data partitions, retaining the dataset's test partition and reserving approximately 20\%
of the training data for validation while preserving class proportions. Exact duplicate content will be checked and handled to prevent overlap across partitions.
Preprocessing will produce three-channel, \(224 \times 224\) images with normalization appropriate to each pretrained checkpoint. Only training images will 
receive augmentation through a common policy of mild rotations, translations, and scaling.

Each architecture will receive the same number of hyperparameter trials and follow a common epoch cap and early-stopping rule. Hyperparameters and checkpoints 
will be selected using validation macro-F1. The final configurations will be evaluated across three fixed random seeds, with test evaluation conducted after model
selection is complete. 

\section*{Evaluation Metric}
Model performance will be evaluated on the held-out test set after hyperparameters and training checkpoints have been selected using the
validation set. Macro-averaged F1 score (macro-F1) will serve as the primary evaluation metric since it combines precision and recall while
assigning equal importance to all four classes. It is defined as
\[
\mathrm{Macro\text{-}F1}
=\frac{1}{4}\sum_{c=1}^{4}
\frac{2TP_c}{2TP_c+FP_c+FN_c},
\]
where \(TP_c\), \(FP_c\), and \(FN_c\) denote the numbers of true positives, false positives, and false negatives for class \(c\), respectively.

We will also report overall accuracy and class-specific precision and recall. Confusion matrices normalized by true class will illustrate misclassification 
patterns. We will summarize classification metrics as the mean and standard deviation across three fixed training seeds using identical data partitions, 
reflecting variation due to training randomness.

Computational efficiency will be assessed using total parameter count, average training time per epoch, and average inference latency per image. Timing 
comparisons will use the same hardware, software environment, input resolution, and numerical precision, with a common training batch size and an inference batch 
size of one. Inference latency will measure model forward passes over repeated trials after warm-up, using GPU synchronization and excluding data loading and GPU 
transfers.

\clearpage
\begingroup
\renewcommand{\bibsection}{\section*{References}}
\begin{thebibliography}{}

\bibitem[He et~al.(2016)]{he2016}
He, K., Zhang, X., Ren, S., \& Sun, J. (2016). Deep residual learning for image recognition.
In 2016 IEEE Conference on Computer Vision and Pattern Recognition (CVPR) (pp. 770--778).
IEEE. \url{https://doi.org/10.1109/CVPR.2016.90}

\bibitem[Liu et~al.(2021)]{liu2021}
Liu, Z., Lin, Y., Cao, Y., Hu, H., Wei, Y., Zhang, Z., Lin, S., \& Guo, B. (2021).
Swin Transformer: Hierarchical vision transformer using shifted windows.
In 2021 IEEE/CVF International Conference on Computer Vision (ICCV) (pp. 9992--10002).
IEEE. \url{https://doi.org/10.1109/ICCV48922.2021.00986}

\bibitem[Mzoughi et~al.(2020)]{mzoughi2020}
Mzoughi, H., Njeh, I., Wali, A., Slima, M. B., BenHamida, A., Mhiri, C., \& Mahfoudhe, K. B. (2020).
Deep multi-scale 3D convolutional neural network (CNN) for MRI gliomas brain tumor classification.
Journal of Digital Imaging, 33(4), 903--915.
\url{https://doi.org/10.1007/s10278-020-00347-9}

\bibitem[Tan and Le(2019)]{tan2019}
Tan, M., \& Le, Q. V. (2019). EfficientNet: Rethinking model scaling for convolutional neural networks.
In K. Chaudhuri \& R. Salakhutdinov (Eds.), Proceedings of the 36th International Conference on Machine Learning
(Vol. 97, pp. 6105--6114). PMLR. \url{https://proceedings.mlr.press/v97/tan19a.html}

\bibitem[Touvron et~al.(2021)]{touvron2021}
Touvron, H., Cord, M., Douze, M., Massa, F., Sablayrolles, A., \& J\'egou, H. (2021).
Training data-efficient image transformers \& distillation through attention.
In M. Meila \& T. Zhang (Eds.), Proceedings of the 38th International Conference on Machine Learning
(Vol. 139, pp. 10347--10357). PMLR. \url{https://proceedings.mlr.press/v139/touvron21a.html}

\end{thebibliography}
\endgroup

\end{document}
